(1) and (2) are \cite[Lem.~7.1]{ButlerT5} read in the direction of construction rather than of analysis: that lemma computes $\Delta = \rho\otimes\bar\rho^{-1}$ and $J|ij\rangle = |ji\rangle$ for exactly this $\Omega$, and establishes that $\Omega$ is cyclic and separating; the identification $\pi(\Aa)' = J\pi(\Aa)J$ is Tomita's theorem, which in this finite-dimensional case is the content of that lemma's polar decomposition. (3) If $|\Psi\rangle \in \C^2\otimes\mathcal{K}$ purifies $\rho$ then its Schmidt coefficients are the $\lambda_i$, so $|\Psi\rangle = (1\otimes W)\Omega$ for an isometry $W$ carrying one Schmidt basis to the other; two purifications on the same purifying space differ by a unitary there. (4) Substituting $\rho = \operatorname{diag}(1, e^{-\omega_k/T})/Z$ gives $\lambda_1/\lambda_0 = e^{-\omega_k/2T} = \tan\theta$, which is the state of \cite[Lem.~7.2]{ButlerT5}.
